\section{Аналитическая часть}\label{sec:analysis}

\subsection{Платформа STaaS и S3-хранилища}\label{subsec:staas}

\gls{s3}~--- это HTTP-протокол для доступа к объектному хранилищу.
Данные в нём представлены неструктурированными объектами с уникальным ключом, бинарным содержимым и метаданными.
Объекты группируются в бакеты~--- именованные контейнеры, выступающие единицей изоляции и управления доступом.
Протокол был предложен Amazon в 2006 году как API к сервису Simple Storage Service~\cite{aws-s3} и быстро стал отраслевым стандартом для объектного хранения.
Сегодня его поддерживает большинство облачных провайдеров и значительная часть открытых решений, в первую очередь Ceph~\cite{weil2006ceph}.

Унификация на уровне протокола позволяет применять одинаковые клиентские библиотеки и инструменты к разнородным хранилищам, а компаниям~--- строить поверх внутренние платформы, не привязываясь к конкретной реализации.
В <<Авито>> такой платформой выступает \gls{staas}.
Она прячет от потребителя детали конкретных кластеров и предоставляет единый API для управления жизненным циклом бакета~--- от его создания и выдачи прав до изменения политик и удаления.
Под капотом \gls{staas} управляет десятками \gls{s3}-кластеров, на которых хранятся петабайтные объёмы данных, разложенные по бакетам.

Управляющим сервисом платформы выступает \textit{service-staas}: через него проходят все операции над бакетами, он хранит метаданные, согласует выдачу доступов в \gls{vault} и поддерживает в актуальном состоянии записи в \gls{consul}~\cite{consul-docs}.
Именно через \gls{consul}-DNS клиент обращается к хранилищу: имя бакета резолвится в адрес конкретного кластера, и далее запрос обрабатывается S3-шлюзом данного кластера.
Такая схема позволяет менять расположение бакета без участия клиента~--- достаточно обновить DNS-запись на стороне платформы.

При этом каждый бакет физически принадлежит ровно одному кластеру.
По мере роста платформы появляется необходимость перемещать бакет между кластерами, сохраняя для клиента ту же точку входа и тот же набор объектов.
Эту задачу и решает механизм миграции, которому посвящены последующие подразделы.

\subsection{Анализ текущего процесса ручной миграции}\label{subsec:manual-migration}

Миграция бакета инициируется по одной из причин: плановое обновление железа, вывод устаревшего кластера или разбиение слишком крупного.
Прежде чем приступить, администратор связывается с владельцем бакета, согласует окно работ и фиксирует за собой ответственность за корректное завершение операции.
С этого момента бакет фактически переходит в режим управляемого переноса, хотя в метаданных самой платформы это никак не отражается~--- сопровождение остаётся <<в голове>> у инженера.

Сама процедура складывается из следующих шагов:
\begin{enumerate}
    \item \textbf{создание целевого бакета} на нужном кластере и выдача доступов, идентичных исходному;
    \item \textbf{запуск \texttt{rclone sync}}~\cite{rclone-docs} из исходного бакета в целевой; для крупных объёмов копирование занимает часы или сутки;
    \item \textbf{контроль прогресса} по логам \gls{rclone}; при сетевых ошибках или сбоях процесс перезапускается администратором вручную;
    \item \textbf{остановка записи} в исходный бакет по согласованию с владельцем, чтобы исключить расхождение данных;
    \item \textbf{финальная досинхронизация} для переноса объектов, появившихся за время основного копирования;
    \item \textbf{переключение DNS-записи} в \gls{consul} на целевой кластер, после чего клиентский трафик начинает идти в новый бакет;
    \item \textbf{сверка и удаление}~--- выборочная проверка содержимого и метаданных, затем планирование исходного бакета на удаление.
\end{enumerate}

Каждый из шагов несёт собственный риск:
\begin{itemize}
    \item \textbf{подготовка доступов и политик} повторяет конфигурацию вручную, и ошибка в одном из правил незаметна до момента, когда клиент столкнётся с отказом на чтение или запись;
    \item \textbf{длительная фаза копирования} зависит от стабильности сети и нагрузки на исходный кластер. Механизма возобновления после сбоя нет, поэтому администратор вынужден держать процесс под наблюдением и при необходимости перезапускать с нуля;
    \item \textbf{окно остановки записи} равно времени финальной досинхронизации. \texttt{rclone sync} не сохраняет состояние между запусками и каждый раз листит весь бакет, поэтому для бакетов с большим числом объектов окно растягивается до заметного простоя на запись;
    \item \textbf{переключение DNS} ошибочной записью приводит к недоступности бакета на стороне клиента, а \textbf{удаление исходного бакета} необратимо~--- любая ошибка в этих шагах оборачивается простоем или потерей данных.
\end{itemize}

Совокупная стоимость миграции складывается не столько из машинного времени, сколько из подготовки, координации и контроля.
При десятках кластеров и регулярных переездах такой подход перестаёт масштабироваться: технически запускать несколько миграций параллельно ничто не мешает, но без статусной модели и наблюдаемости весь контекст приходится держать в голове инженера, что быстро становится источником ошибок.

\subsection{Обзор существующих решений}\label{subsec:rclone-review}

Существующие практики переноса данных между S3-совместимыми хранилищами условно делятся на три класса:
\begin{itemize}
    \item \textbf{утилиты синхронизации уровня объекта} (\texttt{rclone sync}, \texttt{aws s3 sync}~\cite{aws-s3} и аналоги)~--- зрелые CLI-инструменты, рассчитанные на ручной запуск администратором, а не на встраивание в сервис. Они хорошо копируют содержимое бакета, но не управляют жизненным циклом миграции (статусы, DNS, доступы, отчётность). Отсутствие объектных блокировок не позволяет согласовать их работу с клиентским трафиком в бесшовном режиме без форка большой кодовой базы;
    \item \textbf{встроенные механизмы репликации хранилищ} (кластерная и межкластерная репликация Ceph и аналогов)~--- настраиваются на уровне самого хранилища, мимо платформы, и обычно жёстко привязаны к конкретному типу backend; кросс-репликация между разнотипными S3-хранилищами штатно не предусмотрена;
    \item \textbf{управляемые облачные сервисы миграции}~--- удобны для публичных провайдеров, но непригодны во внутреннем контуре закрытой инфраструктуры: они не интегрируются со статусной моделью бакета и системой авторизации \gls{staas} и не дают контроля над объектными блокировками, необходимыми в бесшовном режиме.
\end{itemize}

Поэтому копирование логично вынести за сменный интерфейс движка, а оркестрацию миграции (статусы, доступы, DNS, наблюдаемость) держать внутри платформы.

\begin{table}[H]
\centering
\small
\renewcommand{\arraystretch}{1.3}
\begin{tabular}{|>{\columncolor{gray!15}\bfseries\raggedright\arraybackslash}p{3.4cm}|p{3.0cm}|p{3.0cm}|p{3.0cm}|}
\hline
\rowcolor{gray!15} \textbf{Критерий} & \textbf{Утилиты sync} & \textbf{Нативная репликация} & \textbf{Облачный сервис} \\
\hline
Перенос между разными типами S3 & $\sim$ через адаптеры & $-$ один backend & $\sim$ зависит от провайдера \\
\hline
Жизненный цикл бакета в \gls{staas} & $-$ только данные & $-$ мимо платформы & $\sim$ ограничен API \\
\hline
Прозрачность эксплуатации & $+$ при локальном контроле & $\sim$ логика в backend & $-$ зависит от SLA \\
\hline
\end{tabular}
\caption{Сравнение классов решений для миграции бакетов в \gls{staas}}
\label{tab:solution-comparison}
\end{table}

Прямого аналога автоматизатора миграций для S3-платформы с описанными требованиями в открытом доступе нет.
Однако задача переноса данных между S3-хранилищами сама по себе решается давно.

\subsection{Подходы к автоматизации}\label{subsec:approaches}

Из анализа ручного процесса (подраздел~\ref{subsec:manual-migration}) видно, что автоматизация сводится к двум задачам: копирование данных без участия инженера и переключение трафика.
Главный архитектурный вопрос~--- что делать с записью, идущей в исходный бакет, пока миграция не завершилась.
В зависимости от ответа складываются два подхода:
\begin{itemize}
    \item допустить короткий простой на запись и завершить миграцию финальной досинхронизацией;
    \item проксировать клиентский трафик через специальный компонент, который во время миграции направляет запросы между исходным и целевым бакетом, не прерывая работу клиента.
\end{itemize}

\subsubsection{Миграция с даунтаймом на запись}\label{subsubsec:approach-downtime}

В этом подходе вся механика ручной миграции сохраняется, но управление переходит к оркестратору внутри платформы: он создаёт целевой бакет, запускает копирование, отслеживает прогресс и доводит миграцию до конца.

По сути это тот же ручной процесс, перенесённый на автоматику: добавляется один оркестратор, всё остальное остаётся прежним.

\subsubsection{Бесшовная миграция}\label{subsubsec:approach-seamless}

Для активно пишущих сервисов простой на запись неприемлем.
Чтобы его убрать, между клиентом и бакетом ставится прокси \textit{s3-shifter}: он направляет запросы в исходный или целевой бакет, а мигратор фоном переносит объекты.

Сложность в том, что параллельная клиентская запись и фоновое копирование могут перезаписать друг друга, а готовые инструменты такой ситуации не предусматривают.
Если в подходе с даунтаймом запись на исходный бакет полностью останавливается извне, то здесь нужна точечная синхронизация~--- на уровне отдельных объектов, к которым одновременно обращаются мигратор и клиент.
Поэтому мигратор приходится разрабатывать самостоятельно.

\subsection{Функциональные требования}\label{subsec:func-req}

Администратор может:
\begin{itemize}
    \item запустить миграцию бакета на указанный кластер;
    \item отследить ход миграции и при необходимости отменить её, вернув бакет в исходное состояние;
    \item мигрировать бакет в любое S3-совместимое хранилище.
\end{itemize}

Со стороны владельца бакета:
\begin{itemize}
    \item доступ к бакету сохраняется по тем же реквизитам (имя, секреты), что и до миграции;
    \item ни один объект не теряется, а обновления клиента, сделанные во время миграции, не затираются переносимыми объектами;
    \item для бесшовной миграции~--- запись и чтение не прерываются на всё время переноса.
\end{itemize}

\subsection{Нефункциональные требования}\label{subsec:nonfunc-req}

\begin{itemize}
    \item миграция переживает отказ одного дата-центра без потери прогресса;
    \item возобновление миграции после сбоя без потери прогресса;
    \item во время миграции продолжают соблюдаться действующие \gls{slo} платформы (95\% PUT $\le$ 950\,мс, 98\% GET $\le$ 50\,мс);
    \item наблюдаемость через логи, метрики и трассировки;
    \item TLS на всех каналах передачи данных;
    \item подключение нового типа S3-хранилища не требует переработки ядра мигратора.
\end{itemize}

\subsection{Выводы по главе}\label{subsec:ch1-summary}

В главе разобран ручной процесс миграции бакета на платформе \gls{staas} и показано, что он плохо масштабируется при десятках кластеров: основная стоимость складывается из подготовки, координации и контроля.

Обзор существующих решений показал, что прямого аналога автоматизатора миграций для S3-платформы нет.

Из анализа сформулированы функциональные и нефункциональные требования к автоматизации, состоящей из двух этапов.
